\scanpage{347}
\mspage{12\deleted{\unsure{9}}}

D’autre part, la vérification de
\[
b_0\amalg_{a_0}c_0\unsure{=|\simeq}e_Z\qquad\text{dans }Z^{\wedge}
\]
est immédiate.

On va utiliser cette flèche
\[
\unsure{\Upsilon|Y}\longrightarrow Z
\]
pour construire des exemples intéressants (et ``\unsure{déchirants|délirants}) de sommes amalgamées
\diagram{\begin{tikzcd}[ampersand replacement=\&]
U\arrow[r,hook,"i"]\arrow[d,"p"'] \& X\arrow[d,"q"]\\
U'\arrow[r,hook,"i'"] \& X',
\end{tikzcd}}
avec $i$ \deleted{\illegible{1}} imm. ouverte, \deleted{\illegible{1}} telles qu’on ait même
\[
i:U\longrightarrow\overline U\qquad\text{est asphérique, donc }\in W,
\]
\deleted{\unsure{sans}} \inserted{alors} que pourtant on \deleted{ait} aura
\[
i':U'\unsure{\text{---}|\longrightarrow}\overline{U'}\qquad\notin W,
\]
et ce qui mieux, (ou pis) est,
\[
\deleted{\unsure{q}}p\in W\qquad\text{mais}\qquad\deleted{\unsure{p}}q\notin W.
\]
(et éventuellement, en $=$ temps, $i\in W$ mais $i'\notin W$). De plus (contrairement aux constructions des \unsure{prop.|pp.} 10, 11 (p. 101-107), les catégories envisagées \deleted{seront} \inserted{pourront toutes être} finies. Je suis particulièrement intéressé par le cas où $X$ est un ``\unsure{carré}''
\[
X=\Phi_n\times\Phi_m\qquad(n,m\geq 2)
\]
et où $U$ est son ``bord''
\[
U=\dot X\overset{\mathrm{def}}{=}(\dot\Phi_n\times\Phi_m)\cup(\Phi_n\times\dot\Phi_m).
\]
\deleted{Nous \unsure{reviendrons|avions vu}, sur \unsure{ce|le} cas de \unsure{suite|tête}, et \unsure{précisons|préciserons} les}

% archive sheet number: 347
